\documentclass[11pt]{article}

\usepackage[margin=1in]{geometry}
\usepackage{hyperref}
\usepackage{longtable}

\title{Filesystem-Gated Multi-Agent Workflow for Scientific Software Development}
\author{AI Supervisor/Worker Workflow}
\date{\today}

\renewcommand{\path}[1]{\texttt{#1}}

\begin{document}
\maketitle

\begin{abstract}
This document describes a reusable AI-assisted coding workflow for scientific
software projects.  A milestone-gated supervisor decomposes large goals into
small jobs, one or more implementation workers execute those jobs, and optional
read-only reviewers inspect each worker attempt before acceptance.  The default
deployment uses Codex CLI for every role; Cursor Agent remains an explicit
compatibility option.  All
coordination happens through durable files in the project repository, which
makes state auditable, restartable, and independent of direct agent-to-agent
process control.
\end{abstract}

\section{Design Goals}

The workflow is designed for scientific codes where correctness, reproducibility,
and reviewable history matter more than maximizing autonomous throughput.  The
main goals are:

\begin{itemize}
  \item decompose large milestones into small, closed-form worker jobs;
  \item require tests and scientific assumptions for numerical work;
  \item preserve reviewable Git history and attempt documentation;
  \item keep human intervention at milestone gates rather than every job;
  \item recover from crashes, timeouts, stale locks, and machine restarts;
  \item make reviewer coverage and blocking decisions machine-checkable;
  \item collect metrics for later analysis of agent reliability and cost.
\end{itemize}

\section{Repository Layout}

\begin{longtable}{p{0.32\linewidth}p{0.62\linewidth}}
\textbf{Path} & \textbf{Purpose}\\
\hline
\path{AGENTS.md} & Role, scientific coding, Git hygiene, and workflow-evolution rules.\\
\path{.ai/supervisor/} & Durable supervisor memory: design prompt, project brief, roadmap, ledger, policies, and human gates.\\
\path{.ai/jobs/JNNNN/} & One worker job: task, status, reports, logs, patches, reviewer reports, and feedback.\\
\path{.ai/commit\_docs/} & Markdown documentation for each worker attempt and its commit range.\\
\path{.ai/metrics/} & Structured run metrics and workflow events.\\
\path{.worktrees/JNNNN/} & Isolated Git worktree for a worker job branch.\\
\path{agent\_wrappers/} & Role-neutral wrapper registry with executable, model, and recommended-role metadata.\\
\path{gui/} & Local dashboard for launching/stopping loops, selecting wrappers/models, reviewing gates, and inspecting logs.\\
\path{skills/} & Project skills plus installed reusable workflow skills visible to Codex and Cursor prompts.\\
\path{scripts/} & Dependency-light shell and Python helpers for dispatch, review, metrics, GUI, and recovery.\\
\end{longtable}

\section{Roles}

\subsection{Supervisor}

The supervisor owns planning, dispatch, review, acceptance, rejection, and ledger
maintenance.  In the default configuration it is launched through Codex CLI, but
the dashboard can also launch it through an explicit compatibility wrapper or future
wrapper.  The supervisor reads the project design and current state, creates one
small job at a time, reviews the worker report, tests, diff, commit docs, and
reviewer reports, then integrates or rejects the job.  Supervisor-owned planning
files include \path{roadmap.md}, \path{project\_brief.md}, and \path{ledger.md};
workers may suggest changes to those files but should not own them.

\subsection{Worker}

The worker implements exactly one assigned job in an isolated worktree.  In the
default configuration it is launched through Codex CLI.  Cursor Agent can still be
selected explicitly as a compatibility wrapper.  The worker runs the requested validation,
keeps changes in meaningful commits where possible, and writes a concise report.
Each report includes workflow friction and skill suggestions so repeated
procedural failures can be converted into templates, scripts, protocols, or
skills.

\subsection{Reviewers}

Two optional read-only reviewers run after the worker and tests finish.  Each
reviewer has an independent wrapper and model.  Reviewer A focuses on
scientific, mathematical, numerical, and algorithmic robustness.  Reviewer B
focuses on code quality, build behavior, portability, tests, maintainability,
and Git hygiene.  Reviewers inspect the actual diff rather than relying on the
worker report.

\section{Agent Wrappers and Model Selection}

Agent wrappers decouple workflow roles from a specific command-line product.
The registry in \path{agent\_wrappers/} declares wrapper IDs, labels,
executables, supported roles, recommended roles, default models per role, and
the models that should appear in the dashboard dropdowns.  The built-in wrappers
are role-neutral:

\begin{longtable}{p{0.20\linewidth}p{0.34\linewidth}p{0.36\linewidth}}
\textbf{Wrapper} & \textbf{Recommended roles} & \textbf{Default model behavior}\\
\hline
\path{codex} & Every role & GPT-5.6 Sol for implementation, review, supervision, and modulation; GPT-5.6 Terra for chat.  Narrow routing and classification agents use GPT-5.6 Luna.\\
\path{cursor-agent} & Compatibility only & Retains legacy Cursor model IDs for explicit opt-in; it is not an active default.\\
\end{longtable}

Compatibility wrappers can still be selected explicitly when desired.  Model dropdowns
are populated from the selected wrapper metadata,
which avoids mixing Codex model names with Cursor model names unless a wrapper
explicitly advertises them.

New wrappers, such as a future Claude Code integration, should be added by
creating \path{agent\_wrappers/<wrapper-id>/wrapper.json} and either a command
template or a built-in runner in \path{scripts/agent\_wrapper.py}.  The workflow
loops call \path{scripts/agent\_wrapper.py run}; they do not hard-code a vendor
command except inside the wrapper runner.

\section{Dashboard}

The local dashboard is served by \path{scripts/workflow\_gui.py}.  It provides:

\begin{itemize}
  \item launch and stop controls for worker and supervisor loops;
  \item wrapper and model dropdowns for worker, reviewer A, reviewer B, and supervisor;
  \item current run status, job state, milestone criteria, project tree, and ledger views;
  \item bounded real-time log panels for worker, supervisor, reviewer, and GUI output;
  \item a human milestone-review form with yes/no checklist answers, comments,
        an approval comment, and a major-structural-change override;
  \item workflow-aware chat panels for asking questions before submitting human review.
\end{itemize}

The AIFLOW V2 GUI asks the operating system for a free loopback port when no explicit
\path{--port} is given.  It writes the URL and mutation token to a private,
checkout-scoped \path{ENDPOINT.json}; the token is not printed and the file is removed
when the server exits.  This lets multiple project dashboards run at the same time
without manual port bookkeeping.

When a human gate exists, the worker loop pauses and the supervisor loop exits
with \path{HUMAN\_REVIEW\_REQUIRED}.  This is intentional: the dashboard remains
available for the human review, and after the review is submitted the workflow
can restart the loops automatically according to the project configuration.

\section{Overall Control Flow}

\begin{verbatim}
Human design prompt
        |
        v
Supervisor extracts brief, roadmap, and current milestone
        |
        v
Create exactly one small worker job in .ai/jobs/JNNNN/
        |
        v
Selected worker wrapper implements in .worktrees/JNNNN
        |
        v
Validation, commit docs, changed_files, diffstat, report
        |
        v
Selected reviewer wrappers inspect the actual diff in parallel
        |
        v
Supervisor reviews worker artifacts plus reviewer artifacts
        |
        +-- reject --> feedback.md --> worker retry
        |
        +-- accept --> integrate branch, update ledger
                         |
                         +-- milestone incomplete --> next small job
                         |
                         +-- milestone complete/blocked --> human review gate
\end{verbatim}

\section{Job State Machine}

Each job has a \path{status.json} with immutable \path{base\_sha},
human-readable \path{base\_ref}, the job branch, attempt number, and current
state.  The immutable base commit is the review boundary for all diffs.

\begin{verbatim}
queued/rejected
      |
      v
running -- deterministic or worker failure --> blocked
      |
      v
reviewing -- reviewer timeout/failure/block --> review_failed or review_timeout
      |
      v
ready_for_review -- supervisor rejects --> rejected
      |
      +-- supervisor accepts --> accepted

Any nonterminal job can become cancelled or superseded by explicit decision.
Terminal jobs are not retried by the worker loop.
\end{verbatim}

\section{Worker Attempt Protocol}

The worker loop performs deterministic steps before and after invoking the selected agent:

\begin{enumerate}
  \item Acquire a loop-level singleton lock and a job lock with owner PID.
  \item Create or reuse \path{.worktrees/JNNNN} on branch \path{ai/JNNNN}.
  \item Resolve and persist \path{base\_sha}; never use a moving ref for review.
  \item Run preflight checks: branch, worktree cleanliness, base ancestry, and
        project-configured submodule requirements.
  \item Build a prompt containing only the job task, feedback, visible skills,
        validation command, and report contract.
  \item Invoke the selected worker wrapper with configurable model, timeout, and
        extra local flags.
  \item Commit remaining changes if the selected agent did not create commits.
  \item Run the job validation command, optionally bounded by \path{TEST\_TIMEOUT}.
  \item Record raw and filtered post-test dirty status.  Allowed generated files
        may be declared in \path{allowed\_artifacts.txt}; undeclared dirty files
        block the attempt.
  \item Generate \path{changed\_files}, \path{diffstat}, \path{diff.patch},
        \path{worker\_handoff.attempt-N.json}, \path{report.md}, commit docs,
        consistency checks, metrics, and events.
\end{enumerate}

The raw worker transcript is retained as audit evidence only.  The worker loop
extracts a structured handoff from the final worker report and uses canonical
workflow facts plus that handoff to render \path{report.md} and commit
documentation.  Raw transcript files such as \path{cursor\_final.attempt-N.md}
and \path{cursor\_stream.attempt-N.jsonl} are not embedded into canonical review
artifacts.

\section{Reviewer Protocol}

Reviewers run in parallel through the selected reviewer wrappers.  Each report
must include a fenced YAML block:

\begin{verbatim}
diff_coverage:
  full_diff_reviewed: true
  files_reviewed:
    - path/from/changed_files
  unreviewed_files: []
review_decision:
  recommendation: accept
  blocks_acceptance: false
  blocking_reasons: []
\end{verbatim}

The worker loop checks coverage with \path{check\_reviewer\_coverage.py} and
parses reviewer decisions with \path{analyze\_reviewer\_reports.py}.  Reviewer
timeouts or blocking recommendations prevent the job from becoming ready for
supervisor acceptance.

\section{Supervisor Acceptance Protocol}

The supervisor reviews the canonical artifacts rather than relying on a single
report.  It should inspect:

\begin{itemize}
  \item \path{status.json}, including \path{base\_sha}, attempts, tests, and reviewer fields;
  \item \path{worker\_handoff.attempt-N.json}, \path{report.md}, \path{test.attempt-N.log}, \path{diffstat}, and selected patch sections;
  \item \path{changed\_files.attempt-N.txt} and reviewer coverage;
  \item reviewer A and B reports and reviewer-decision analysis;
  \item commit documentation and the actual branch diff \path{base\_sha..HEAD};
  \item attempt-consistency and post-test-dirty artifacts.
\end{itemize}

Before integration, \path{integrate\_job.py} verifies state, tests, reviewer
completion, post-test cleanliness, attempt consistency, branch existence, and
base ancestry.  Integration is explicit and should be done from a clean main
worktree.

\section{Human Gates}

Human review is reserved for milestone boundaries, major structural changes, and
decisions the supervisor cannot safely make.  The dashboard cycles through every
checklist item, records yes/no answers and comments, and supports an optional
approval comment.  Ordinary failed checklist items are routed through
\path{HUMAN\_REVIEW\_ACTION\_REQUESTED.md}; the supervisor classifies them and
creates a small revision job, updates planning records, or opens a clarification
gate.  Major structural requests supersede the checklist outcome and are routed
through \path{STRUCTURAL\_CHANGE\_REQUESTED.md}.  The supervisor must update the
milestone plan first, then return to the human with a revised review gate before
implementation continues.

\section{Failure Handling}

\begin{longtable}{p{0.30\linewidth}p{0.64\linewidth}}
\textbf{Failure mode} & \textbf{Handling}\\
\hline
Worker timeout & Record timeout, optionally auto-resume if configured, otherwise block.\\
Worker crash & Auto-relaunch limited attempts for non-hard failures; record event.\\
Reviewer timeout/crash & Retry reviewers, then set \path{review\_timeout} or \path{review\_failed}.\\
Reviewer blocks acceptance & Keep job out of \path{ready\_for\_review}; supervisor decides reject, waive, or rerun.\\
Post-test dirty files & Block unless all paths match \path{allowed\_artifacts.txt}.\\
Stale running job & Recover only if lock owner PID is gone; mark \path{implemented} if artifacts exist.\\
Supervisor crash & Create \path{SUPERVISOR\_ACTION\_REQUIRED.md}, pause worker dispatch, and avoid a human milestone gate unless a real human decision is needed.\\
Superseded work & Use terminal state \path{superseded} so the worker loop does not retry it.\\
\end{longtable}

\section{Workflow Evolution}

Workers report friction and skill suggestions.  Reviewers evaluate those
suggestions.  The supervisor deduplicates them against existing skills,
templates, scripts, and protocols, then chooses the smallest durable fix:
ledger note, documentation, checklist, template, protocol, script, project skill,
or reusable workflow skill.  Repeated failure artifacts can also trigger
supervisor-originated improvements, even when workers do not recognize the
pattern.

\section{Metrics and Paper Reproducibility}

Worker and reviewer runs write per-agent metrics next to job artifacts, and
supervisor runs write metrics under \path{.ai/metrics/supervisor/}.  All records
are appended to \path{.ai/metrics/runs.jsonl}.  Workflow events such as failures,
reviewer blocks, wrapper crashes, and human interventions are appended to
\path{.ai/metrics/events.jsonl}.  These records can support a later empirical
study of agent reliability, failure modes, human-intervention frequency, token
use, wall-clock time, wrapper choice, model choice, and scientific-code review
outcomes.

\section{Minimal Commands}

\begin{verbatim}
python3 scripts/workflow_gui.py --host 127.0.0.1 --project-root .

WORKER_AGENT_WRAPPER=codex \
WORKER_MODEL=gpt-5.6-sol \
REVIEWER_A_AGENT_WRAPPER=codex \
REVIEWER_A_MODEL=gpt-5.6-sol \
REVIEWER_B_AGENT_WRAPPER=codex \
REVIEWER_B_MODEL=gpt-5.6-sol \
./scripts/worker_loop.sh

SUPERVISOR_AGENT_WRAPPER=codex \
CODEX_MODEL=gpt-5.6-sol \
CODEX_REASONING_EFFORT=high \
./scripts/supervisor_loop.sh

python3 scripts/agent_wrapper.py list --json
python3 scripts/summarize_jobs.py
python3 scripts/summarize_agent_metrics.py
\end{verbatim}

\end{document}
